\chapter*{Abstract}
 
Delays are inherent to most dynamical systems. 
Besides shifting the process in time, they can significantly affect their performance. 
For this reason, it is usually valuable to study the delay and account for it.
Because they are dynamical systems, it is of no surprise that sequential decision-making problems such as \glspl{mdp} can also be affected by delays.
These processes are the foundational framework of \gls{rl}, a paradigm whose goal is to create artificial agents capable of learning to maximise their utility by interacting with their environment.

\glspl{rl} has achieved strong, sometimes astonishing, empirical results, but delays are seldom explicitly accounted for. 
The understanding of the impact of delay on the \gls{mdp} is limited.
In this dissertation, we propose to study the delay in the agent's observation of the state of the environment or in the execution of the agent's actions.
We will repeatedly change our point of view on the problem to reveal some of its structure and peculiarities.
A wide spectrum of delays will be considered, and potential solutions will be presented.
This dissertation also aims to draw links between celebrated frameworks of the \gls{rl} literature and the one of delays.
We will therefore focus on the following four points of view. 

At first, we consider constant delays. Taking a psychology-inspired approach, we study the impact of predicting the near future in order to estimate the impact of the agent's actions.
We will highlight how this approach relates to models from the \gls{rl} literature. 
It will also be the occasion to formally demonstrate a seemingly evident fact: longer delays result in lower performances. 
The experimental analysis will conclude by showing the validity of the approach.

As a second point of view, we will consider the simple approach of imitating an undelayed expert behaviour in the delayed environment.
The delay will remain constant at first, but we will extend our study to more exotic types of delay, such as stochastic ones.
Although simple, we demonstrate the great theoretical guarantees and empirical results of the approach.

For our third perspective, again on constant delay, we consider adopting a non-stationary memoryless behaviour.
Although it seemingly ignores the delay, the approach treats the delay's effect as an unobserved variable that guides its non-stationarity.
Building on this idea, we provide a theoretically grounded algorithm for learning such behaviour that we test in realistic scenarios. 

Finally, our last point of view will consider a broader model than that of constant delay, which includes constant delays as a special case. 
This model will enable actions to affect multiple future transitions of the environment.
Its theoretical properties will be examined to understand its specificities.
Based on these properties, some \gls{rl} algorithms will be ruled out, while others will be tested in various empirical studies. 
As a more general model for delays, its understanding has implications for the constant delay frameworks of the previous chapters. 
